\chapter{Ankle Arthrodesis}

\section*{What Is This Surgery?}
Ankle arthrodesis, commonly known as ankle fusion, is a definitive surgical procedure designed to permanently eliminate motion in the tibiotalar joint, which is the primary articulation of the ankle. This operation involves the surgical removal of the remaining articular cartilage from the ends of the distal tibia and the talus bone, followed by their rigid fixation in an optimal functional position, allowing them to grow together into a single, solid bone. The primary goal of ankle arthrodesis is to relieve severe, debilitating pain, correct significant deformity, and restore stability in patients suffering from end-stage ankle arthritis or other conditions that have rendered the joint non-functional and persistently painful. It is considered a highly effective and durable solution for appropriate candidates, providing a stable, pain-free platform for ambulation, albeit at the cost of joint motion.

\section*{Alternative Names}
\begin{itemize}
  \item Ankle fusion
  \item Tibiotalar arthrodesis
  \item Tibiotalar fusion
  \item Pantalar arthrodesis (if subtalar and/or talonavicular joints are also fused)
  \item Triple arthrodesis (a distinct procedure involving subtalar, talonavicular, and calcaneocuboid joints, sometimes performed concurrently or sequentially)
\end{itemize}

\section*{Relevant Surgical Anatomy}
The tibiotalar joint, the focus of ankle arthrodesis, is a hinge-type synovial joint formed by the articulation of the distal tibia and fibula (forming the ankle mortise) with the trochlea of the talus. The distal tibia contributes the plafond (weight-bearing surface) and the medial malleolus, while the distal fibula forms the lateral malleolus. The talus, a unique bone with no muscular attachments, serves as a critical link between the leg and the foot, transmitting forces from the tibia to the calcaneus. Key ligaments stabilizing this joint include the deltoid ligament medially (tibionavicular, tibiocalcaneal, anterior and posterior tibiotalar parts) and the lateral collateral ligaments (anterior talofibular, posterior talofibular, and calcaneofibular ligaments).

Surrounding the joint are several important neurovascular structures. Anteriorly, the anterior tibial artery, deep peroneal nerve, and extensor tendons (tibialis anterior, extensor hallucis longus, extensor digitorum longus) traverse the ankle. Medially, the posterior tibial artery, tibial nerve, and flexor tendons (tibialis posterior, flexor digitorum longus, flexor hallucis longus) pass behind the medial malleolus. Laterally, the superficial peroneal nerve and sural nerve are vulnerable during surgical approaches. The peroneal artery lies deep to the fibula. Lymphatic drainage of the ankle and foot primarily follows the venous system to popliteal and inguinal lymph nodes. Surgical landmarks include the palpable malleoli, the anterior joint line, and the various tendon sheaths that guide surgical approaches and protect underlying structures. Understanding the precise three-dimensional relationship of these structures is paramount to minimize iatrogenic injury during dissection and fixation.

\section{Surgical Principle \& Rationale}
The fundamental principle of ankle arthrodesis is to convert a painful, unstable, or severely deformed mobile joint into a stable, pain-free, immobile bony union. This is achieved by surgically preparing the opposing bone surfaces of the distal tibia and the superior talus, meticulously removing all remnants of diseased articular cartilage and exposing healthy, bleeding subchondral bone. This exposure of cancellous bone is crucial as it provides the necessary osteogenic cells and vascularity for bone healing. These prepared surfaces are then meticulously aligned in a position that optimizes foot function for walking, typically in slight dorsiflexion (0-5 degrees) and neutral rotation, with minimal valgus or varus, to minimize stress on adjacent joints and facilitate a relatively normal gait pattern.

Once aligned, the bone ends are compressed and held rigidly together using various internal fixation devices, such as screws, plates, or intramedullary nails. This rigid fixation is paramount to prevent micromotion at the fusion site, which can impede bone healing and lead to nonunion. The biological process of osteointegration, where the two bones grow together, is then allowed to occur over several weeks to months. Bone graft, either autograft (from the patient's own body, e.g., iliac crest, distal tibia) or allograft (from a donor), may be utilized to enhance bone healing, particularly in cases of significant bone loss, large defects, or revision surgery where the risk of nonunion is higher. The ultimate technical goal is to achieve a solid, stable bony bridge across the former joint space, thereby eliminating the source of pain and providing a durable foundation for weight-bearing and ambulation.

\section{History \& Surgical Evolution}
The concept of joint fusion for pain relief and stability dates back to ancient times, with evidence found in skeletal remains. However, formal ankle arthrodesis procedures began to emerge in the late 19th century. Early pioneers like Dr. Albert in 1879 and Dr. Ollier in 1881 described techniques for ankle fusion, primarily for tuberculosis of the ankle joint, which was a prevalent cause of joint destruction at the time. These initial methods often involved extensive open approaches and external fixation, with variable success rates due to challenges in achieving rigid fixation and effectively managing postoperative infection.

Throughout the 20th century, surgical techniques evolved significantly. The introduction of internal fixation devices, such as screws and plates, by surgeons like Charnley in the mid-20th century, revolutionized arthrodesis by providing more stable constructs and dramatically improving fusion rates. The development of intramedullary nails in the latter half of the century offered another robust fixation option, particularly for complex cases, revision surgeries, or in patients with poor bone quality. More recently, minimally invasive and arthroscopic techniques have been developed, aiming to reduce soft tissue dissection, decrease postoperative pain, and potentially accelerate recovery, though open techniques remain the gold standard for many indications, especially those with significant deformity or bone loss. The evolution reflects a continuous effort to improve fusion rates, reduce complications, and enhance patient recovery.

\section{Clinical Indications}
Ankle arthrodesis is typically indicated for patients experiencing severe, chronic ankle pain, instability, or deformity that has not responded to comprehensive conservative management.

\begin{itemize}
  \item End-stage primary osteoarthritis of the ankle, characterized by severe cartilage loss, subchondral sclerosis, and bone-on-bone articulation, leading to debilitating pain.
  \item Post-traumatic arthritis, which is the most common cause of ankle arthritis, resulting from previous fractures, severe sprains, or repetitive microtrauma to the joint.
  \item Inflammatory arthropathies, such as rheumatoid arthritis, psoriatic arthritis, or ankylosing spondylitis, when medical management fails to control joint destruction and pain.
  \item Neuropathic arthropathy (Charcot joint), particularly in diabetic patients, causing severe joint destruction, instability, and risk of ulceration or infection.
  \item Severe ankle instability or deformity (e.g., varus or valgus malalignment) that cannot be corrected by soft tissue procedures or osteotomies, leading to functional impairment.
  \item Failed total ankle arthroplasty, where revision to fusion may be necessary due to infection, aseptic loosening of components, persistent pain, or severe periprosthetic osteolysis.
  \item Avascular necrosis of the talus, leading to collapse of the talar dome and subsequent severe arthritis.
  \item Tumors involving the ankle joint requiring resection and subsequent stabilization.
  \item Chronic osteomyelitis of the tibiotalar joint that is refractory to debridement and antibiotic therapy.
\end{itemize}

\section*{Clinical Positioning}
Ankle arthrodesis can serve as both a primary and a secondary surgical treatment, depending on the patient's specific condition, activity level, and comorbidities. It is often considered a primary definitive treatment for end-stage ankle arthritis in younger, more active individuals who place high demands on their ankle, as the fusion provides a durable and stable construct. It is also preferred in patients with conditions that contraindicate total ankle arthroplasty (TAA), such as active infection, severe bone loss, significant deformity, severe neuropathic joint disease (Charcot arthropathy), or poor soft tissue envelope, where the durability and stability of a fusion are preferred over motion preservation.

As a secondary or salvage procedure, ankle arthrodesis is frequently performed following a failed total ankle arthroplasty, where it provides a reliable solution for pain relief and stability after implant failure or infection. It is also a secondary option for complex post-traumatic deformities, chronic infections that have not responded to initial treatments, or in cases of nonunion after previous ankle fracture fixation. While TAA aims to preserve motion, arthrodesis prioritizes pain relief and stability, making it a robust alternative when motion preservation is not feasible, not durable enough for the patient's lifestyle, or not desirable due to underlying pathology. The decision between arthrodesis and arthroplasty is complex and individualized, weighing the benefits of motion preservation against the long-term stability and pain relief offered by fusion.

\section*{Surgical Technique \& Procedure}
Ankle arthrodesis is performed under general anesthesia, often supplemented with a regional nerve block (e.g., popliteal or ankle block) for enhanced postoperative pain control. The patient is typically positioned supine with a tourniquet applied to the thigh, though lateral or prone positions may be used depending on the chosen surgical approach and surgeon preference. Common approaches include the anterior approach (between the tibialis anterior and extensor hallucis longus tendons), a lateral transfibular approach (involving fibular osteotomy), or a posterior approach (for specific deformities or revision cases).

The key operative steps generally include:
\begin{itemize}
  \item \textbf{Incision and Exposure:} A surgical incision is made according to the chosen approach, carefully dissecting through subcutaneous tissues while protecting neurovascular structures. The joint capsule is incised to expose the tibiotalar articulation.
  \item \textbf{Cartilage Resection and Joint Preparation:} All remaining articular cartilage from the distal tibia plafond and the superior talar dome is meticulously removed using osteotomes, burrs, or curettes, exposing healthy, bleeding subchondral bone. Any osteophytes (bone spurs) or sclerotic bone are also resected to ensure optimal bone-to-bone contact. The bone surfaces are contoured to ensure maximal contact and stability, often creating flat or slightly concave/convex surfaces.
  \item \textbf{Deformity Correction and Alignment:} The ankle is carefully positioned to correct any pre-existing deformity and achieve a functional alignment. This typically involves 0-5 degrees of dorsiflexion, neutral rotation, and neutral valgus/varus, which is critical for optimizing gait mechanics and minimizing stress on adjacent joints.
  \item \textbf{Bone Grafting (Optional):} Autogenous bone graft (e.g., from the iliac crest, distal tibia, or calcaneus) or allograft may be packed into any gaps or defects to promote fusion, especially in cases of significant bone loss, large cysts, or when there is a high risk of nonunion.
  \item \textbf{Internal Fixation:} The prepared bone surfaces are compressed and stabilized with rigid internal fixation. This commonly involves multiple large cancellous screws placed across the fusion site in a divergent pattern, or a plate and screw construct applied to the anterior or lateral aspect of the joint, or an intramedullary nail inserted antegrade through the heel into the talus and tibia. The choice of fixation depends on the approach, bone quality, and surgeon preference.
  \item \textbf{Wound Closure:} The surgical site is thoroughly irrigated with saline. The incision is closed in layers, typically with absorbable sutures for deep layers and non-absorbable sutures or staples for the skin. A drain may be placed temporarily to manage postoperative bleeding and reduce hematoma formation. A sterile dressing and a well-padded splint or cast are applied to maintain immobilization.
\end{itemize}

\section*{Preoperative Workup \& Preparation}
Thorough preoperative preparation is essential to optimize outcomes and minimize complications for ankle arthrodesis. This typically involves a multi-faceted approach:

\begin{itemize}
  \item \textbf{Comprehensive Medical Evaluation:} A detailed history and physical examination, including assessment of cardiovascular, pulmonary, renal, and endocrine function. This often requires medical clearance from a primary care physician or cardiologist, especially for patients with significant comorbidities.
  \item \textbf{Laboratory Tests:} Routine blood work, including complete blood count (CBC), electrolyte panel, renal function tests, liver function tests, coagulation studies (PT/INR, PTT), and blood type and screen. Hemoglobin A1c is crucial for diabetic patients to assess glycemic control.
  \item \textbf{Imaging Studies:} Standard weight-bearing ankle X-rays (AP, lateral, mortise views) are essential to assess joint space narrowing, osteophyte formation, and deformity. These are often supplemented with CT scans to provide detailed information on bone quality, subchondral cysts, precise joint space narrowing, and three-dimensional deformity, which aids in surgical planning. MRI may be used to evaluate soft tissues, avascular necrosis, or occult infection.
  \item \textbf{Anesthesia Consultation:} A meeting with the anesthesiologist to discuss anesthesia options (general, regional, or combined), address any patient concerns, and optimize medical conditions for surgery.
  \item \textbf{Medication Review:} Adjustment or cessation of certain medications, such as anticoagulants (e.g., warfarin, DOACs) or anti-inflammatory drugs (NSAIDs), as advised by the surgeon and anesthesiologist, to minimize bleeding risk and potential impairment of bone healing.
  \item \textbf{Smoking Cessation:} Patients are strongly advised to stop smoking several weeks to months before surgery, as nicotine significantly impairs bone healing, increases the risk of nonunion, and compromises wound healing.
  \item \textbf{Nutritional Optimization:} Addressing any nutritional deficiencies, particularly vitamin D and protein, to support bone healing and overall recovery.
  \item \textbf{Preoperative Antibiotics:} Administered intravenously shortly before the incision (typically within 60 minutes) to reduce the risk of surgical site infection.
  \item \textbf{Informed Consent:} A detailed discussion of the procedure, potential risks, expected benefits, and alternative treatments, ensuring the patient fully understands and agrees to the surgery.
  \item \textbf{NPO Status:} Patients must refrain from eating or drinking for a specified period (typically 6-8 hours) before surgery to prevent aspiration during anesthesia.
\end{itemize}

\section*{Contraindications \& Precautions}
\textbf{Absolute Contraindications:}
\begin{itemize}
  \item Active infection in the ankle joint or surrounding soft tissues (e.g., osteomyelitis, cellulitis), as this significantly increases the risk of surgical site infection and nonunion.
  \item Severe peripheral vascular disease with critical limb ischemia, non-healing ulcers, or poor perfusion, which compromises wound healing and the biological potential for fusion.
  \item Uncontrolled diabetes mellitus (e.g., HbA1c > 8-9\%), leading to high infection risk, impaired bone healing, and increased risk of Charcot neuroarthropathy progression.
  \item Severe osteopenia or osteoporosis that precludes stable internal fixation, making it impossible to achieve rigid compression across the fusion site.
  \item Acute Charcot arthropathy with ongoing bone destruction and inflammation, as fusion attempts in this phase often fail.
  \item Patient non-compliance or inability to adhere to strict non-weight-bearing protocols and postoperative care instructions.
\end{itemize}

\textbf{Relative Contraindications and Precautions:}
\begin{itemize}
  \item Heavy smoking, which significantly increases the risk of nonunion (up to 30-50\%) and wound complications. Patients are strongly encouraged to quit.
  \item Poor nutritional status or morbid obesity, which can impair healing, increase surgical risks, and make rehabilitation more challenging.
  \item Significant arthritis in adjacent joints (e.g., subtalar, midfoot, or knee), which may become symptomatic or worsen after ankle fusion due to increased compensatory motion and stress.
  \item Neurological deficits that severely impair ambulation or sensation, making postoperative rehabilitation challenging and increasing the risk of pressure sores.
  \item Chronic regional pain syndrome (CRPS) in the affected limb, as surgery can exacerbate symptoms.
  \item Active substance abuse, which can impact compliance, pain management, and overall recovery.
  \item Immunosuppression (e.g., due to medications or systemic disease), increasing infection risk.
\end{itemize}

\section*{Complications \& Risks}
As with any surgical procedure, ankle arthrodesis carries potential risks, which can be broadly categorized as general surgical risks and procedure-specific complications.

\textbf{General Surgical Risks:}
\begin{itemize}
  \item \textbf{Anesthesia Complications:} Adverse reactions to anesthetic agents, respiratory or cardiac compromise, aspiration, and nerve damage from positioning.
  \item \textbf{Bleeding:} Excessive blood loss during or after surgery, potentially requiring blood transfusion or re-operation for hematoma evacuation.
  \item \textbf{Infection:} Superficial wound infection (cellulitis) or deep surgical site infection (osteomyelitis), which can be severe, requiring prolonged antibiotic therapy, debridement, or even amputation in rare, recalcitrant cases.
  \item \textbf{Deep Vein Thrombosis (DVT) and Pulmonary Embolism (PE):} Blood clots forming in the deep veins of the leg, which can dislodge and travel to the lungs, potentially causing life-threatening respiratory compromise.
  \item \textbf{Nerve or Vascular Injury:} Direct damage to nerves (e.g., superficial peroneal, sural, saphenous) causing numbness, weakness, or chronic neuropathic pain, or injury to blood vessels leading to ischemia or hemorrhage.
\end{itemize}

\textbf{Procedure-Specific Risks:}
\begin{itemize}
  \item \textbf{Nonunion or Delayed Union:} Failure of the bones to fuse completely, or taking longer than expected, requiring further surgery (revision arthrodesis, bone grafting, or hardware exchange). This is the most common complication, with rates ranging from 5\% to 30\%, significantly higher in smokers or diabetics.
  \item \textbf{Malunion:} Fusion of the ankle in an incorrect or non-functional position (e.g., excessive varus/valgus, equinus), leading to altered gait mechanics, pain, and increased stress on adjacent joints. May require corrective osteotomy.
  \item \textbf{Hardware Complications:} Breakage, loosening, migration, or prominence of screws, plates, or nails, necessitating removal or revision surgery.
  \item \textbf{Adjacent Joint Arthritis:} Increased stress and compensatory motion in the subtalar, midfoot, or knee joints, potentially leading to new or accelerated arthritis over time. This is a common long-term sequela.
  \item \textbf{Chronic Pain:} Persistent pain at the fusion site or in adjacent areas, sometimes due to nerve irritation, hardware prominence, or complex regional pain syndrome (CRPS).
  \item \textbf{Wound Healing Problems:} Delayed healing, dehiscence (wound opening), or skin necrosis, especially in patients with poor circulation, diabetes, or those undergoing revision surgery.
  \item \textbf{Limb Length Discrepancy:} Shortening of the affected limb, which may require a shoe lift to equalize leg lengths and prevent compensatory gait abnormalities.
  \item \textbf{Stiffness in Adjacent Joints:} While the ankle joint is fused, other foot joints may also experience some degree of stiffness due to prolonged immobilization or compensatory changes.
  \item \textbf{Pseudarthrosis:} A false joint forming at the fusion site due to nonunion, often presenting with persistent pain and instability.
\end{itemize}

\section*{Postoperative Recovery Timeline}
Immediately after ankle arthrodesis, the patient is transferred to the Post-Anesthesia Care Unit (PACU) for close monitoring of vital signs, pain levels, and neurovascular status of the foot (color, temperature, capillary refill, sensation, motor function). The leg will be immobilized in a well-padded splint or cast, and strict elevation of the limb above heart level is crucial to minimize swelling and reduce pain. Pain management is a priority, often involving a multimodal approach with a combination of oral and intravenous medications, and the regional nerve block typically provides several hours of initial relief. Early mobilization of the patient, usually to a chair, is encouraged to prevent complications like DVT.

Over the first 24-48 hours, the patient typically remains hospitalized. Nurses will continue to monitor the surgical site for signs of bleeding or infection, and administer prophylactic antibiotics and anticoagulants to prevent DVT. Discharge planning begins early, focusing on safe mobility with crutches or a walker, proper cast care, and effective pain control. For the first 6-12 weeks, strict non-weight-bearing is typically required to allow the bones to begin fusing without disruption. The initial splint is often replaced with a non-weight-bearing cast or controlled ankle motion (CAM) boot after the first week or two, once initial wound healing is confirmed. Regular follow-up appointments with the surgeon are scheduled to monitor wound healing and assess fusion progress with serial X-rays. Gradual weight-bearing is initiated only when radiographic evidence of early fusion is present, progressing from partial to full weight-bearing over several weeks, often with the aid of physical therapy to restore gait mechanics, strengthen surrounding muscles, and adapt to the fused ankle. Full bony union can take 3 to 6 months, or even longer in some cases, with complete recovery and adaptation extending up to a year.

\section*{Surgical Findings \& Pathology}
During ankle arthrodesis, the surgeon meticulously documents the intraoperative findings, which typically include the extent of articular cartilage loss, subchondral bone sclerosis, presence and size of osteophytes, joint space narrowing, and any associated synovitis or cyst formation. The degree of deformity (e.g., varus, valgus, equinus) and its reducibility are also noted. Any unexpected findings, such as occult infection or significant bone defects, are addressed and documented. The resected articular cartilage and subchondral bone fragments are routinely sent for pathological examination.

The pathology report on the resected tissues typically confirms the diagnosis of degenerative osteoarthritis, post-traumatic arthritis, or inflammatory arthropathy. Histologically, this involves findings such as fragmented articular cartilage, subchondral bone necrosis or sclerosis, synovial hyperplasia, and inflammatory cell infiltrates depending on the underlying etiology. If bone graft was used, the report may describe its integration. In cases of suspected infection, cultures are taken intraoperatively, and the pathology report may show evidence of acute or chronic inflammation consistent with osteomyelitis. The pathology findings provide definitive confirmation of the disease process and guide any further management if unexpected results are found.

\section*{Expected Postoperative Course}
A normal and successful postoperative course following ankle arthrodesis is characterized by a progressive reduction in pain, leading to comfortable weight-bearing and ambulation once fusion is achieved. Initially, patients will experience surgical pain managed by medication, but the severe, chronic pain present preoperatively should gradually resolve as the joint stabilizes. The timeline for healing involves several phases: initial wound healing (2-4 weeks), early bone healing with callus formation (6-12 weeks, non-weight-bearing), and progressive consolidation of the fusion (3-6 months, gradual weight-bearing). Radiographic evidence of solid bony bridging across the tibiotalar joint is the ultimate confirmation of success. Patients should expect to adapt to the absence of ankle motion, with compensatory movements occurring in the subtalar and midfoot joints, allowing for a relatively normal gait pattern on flat surfaces. The goal is to achieve a stable, pain-free ankle platform that significantly improves the patient's quality of life and functional independence.

\section{Postoperative Care \& Follow-up}
Postoperative care and diligent follow-up are critical to monitor healing, prevent complications, and ensure optimal long-term outcomes after ankle arthrodesis.

\begin{itemize}
  \item \textbf{Regular Clinical Visits:} Typically scheduled at 2 weeks, 6 weeks, 3 months, 6 months, and 1 year post-surgery. These visits allow the surgeon to assess wound healing, pain levels, and functional progress.
  \item \textbf{Serial X-rays:} Performed at each follow-up visit to monitor the progress of bone healing and confirm solid fusion. CT scans may be ordered if fusion is delayed, uncertain, or if nonunion is suspected.
  \item \textbf{Suture/Staple Removal:} Usually occurs at the 2-week post-op visit, provided the wound is healing well.
  \item \textbf{Physical Therapy (PT):} Initiated once partial weight-bearing is allowed (typically 6-12 weeks post-op). PT focuses on gait training, strengthening of surrounding muscles (e.g., quadriceps, hamstrings, hip abductors), balance exercises, and adaptation to the fused ankle.
  \item \textbf{Activity Resumption:} Gradual return to daily activities and low-impact exercises (e.g., swimming, cycling) as fusion progresses and pain subsides. High-impact activities (e.g., running, jumping sports) are generally discouraged long-term to protect adjacent joints.
  \item \textbf{Footwear Modifications:} Patients may require supportive shoes, custom orthotics, or rocker-bottom soles to optimize comfort, gait mechanics, and reduce stress on adjacent joints.
\end{itemize}
Patients should be educated on warning signs to report immediately, such as increasing pain, redness, swelling, warmth, fever, purulent drainage from the wound, or new numbness/weakness in the foot, as these may indicate infection or other complications.

\section*{Epidemiology}
Ankle arthrodesis is a well-established procedure, though less common than hip or knee arthroplasty. The incidence of end-stage ankle arthritis, the primary indication for fusion, is estimated to affect approximately 1\% of the adult population, with post-traumatic arthritis being the leading cause, accounting for over 70\% of cases. This high prevalence of post-traumatic arthritis means ankle arthrodesis is often performed in a younger patient demographic (typically 40-60 years old) compared to total ankle arthroplasty, with a higher prevalence in males due to the higher incidence of traumatic ankle injuries. While precise epidemiological data on arthrodesis rates vary by region, it remains a significant surgical option for end-stage ankle disease. Morbidity rates are primarily driven by nonunion, which can range from 5\% to 30\% depending on patient factors (e.g., smoking, diabetes, bone quality) and surgical technique. Other significant morbidities include wound complications (5-15\%) and adjacent joint arthritis (30-50\% long-term). Mortality directly attributable to ankle arthrodesis is exceedingly low, primarily associated with general anesthesia risks in medically complex patients.

\section*{Outcomes \& Prognosis}
The outcomes of ankle arthrodesis are generally excellent in terms of pain relief and functional stability, making it a highly effective treatment for end-stage ankle arthritis. Success rates for achieving a solid bony fusion range from 80\% to 95\%, with a corresponding high rate of patient satisfaction regarding pain reduction (typically 80-90\% report significant improvement or complete pain relief). While the procedure results in permanent loss of tibiotalar joint motion, most patients adapt well, with compensatory motion occurring in the subtalar and midfoot joints, allowing for a relatively normal gait on flat surfaces. Functional outcomes include improved walking distance, reduced reliance on assistive devices, and enhanced quality of life, often enabling patients to return to light recreational activities.

The long-term prognosis is generally favorable, with the fusion providing a durable solution for pain and stability. However, a significant long-term concern is the development of adjacent joint arthritis, particularly in the subtalar and midfoot joints, due to increased stress and compensatory motion. This can occur in 30-50\% of patients over 10-20 years, potentially requiring further intervention such as subtalar fusion. Prognostic factors significantly influencing fusion rates and long-term outcomes include smoking status (a major negative predictor, increasing nonunion risk by 2-4 times), obesity, diabetes, bone quality, and the surgeon's experience. Achieving optimal alignment and rigid fixation are crucial for a successful and lasting result. While no single landmark clinical trial like APPAC or NASCET exists for ankle arthrodesis, numerous cohort studies and systematic reviews consistently demonstrate its efficacy and durability for appropriate indications.

\section*{References}
\begin{enumerate}
  \item MedlinePlus. Ankle fusion. U.S. National Library of Medicine; 2024. Available from: \url{https://medlineplus.gov/ency/article/007204.htm}
  \item Campbell's Operative Orthopaedics. 14th ed. Philadelphia, PA: Elsevier; 2021.
  \item Rockwood and Green's Fractures in Adults. 9th ed. Philadelphia, PA: Wolters Kluwer; 2020.
  \item American Academy of Orthopaedic Surgeons (AAOS). Ankle Arthrodesis. OrthoInfo. Available from: \url{https://orthoinfo.aaos.org/en/diseases--conditions/ankle-arthrodesis/}
  \item Saltzman CL, et al. Ankle Arthrodesis. J Bone Joint Surg Am. 2005;87(2):401-416.
\end{enumerate}

\clearpage